% ============================================================================
%   Kernel-Level Acceleration: From Valid 2:4 Masks to Sparse Tensor Core Dispatch
%   Standalone draft — to be inserted into final (1).tex
%   Location: Appendix G, after Hardware Profiling Protocol (line ~1415),
%             before Quality Gated Failure Detection Protocol (line ~1417)
% ============================================================================

\subsection{Kernel-Level Acceleration: From Valid 2:4 Masks to Sparse Tensor Core Dispatch}
\label{app:kernel_acceleration}

DST-EDL establishes exact valid block 2:4 feasibility by construction
(Appendix~\ref{app:24_feasibility}). This subsection formalizes the
additional conditions required to translate that structural guarantee into
realized inference acceleration through Sparse Tensor Core dispatch via
cuSPARSELt~\cite{nvidiacusparselt}, TensorRT sparse execution, or PyTorch
semi-structured kernels~\cite{pytorchsemistructured}.

\paragraph{Compressed storage format and metadata encoding.}
Sparse Tensor Cores on Ampere and later architectures operate on a compressed
representation~\cite{mishra2021accelerating}. For each contiguous 4-block
$(w_{4k},w_{4k+1},w_{4k+2},w_{4k+3})$ in a weight row that satisfies the 2:4
constraint, the hardware stores only the two nonzero values and a 2-bit index
that encodes their original positions within the block. Formally, let
$M\in\mathcal{C}_{2:4}$ be the mask produced by Eq.~\ref{eq:mask}. For each
row $i$ and block $k$, define the index set
$I_{i,k}=\{j\in\{0,1,2,3\}:M_{i,4k+j}=1\}$ with $|I_{i,k}|=2$. The
compressed representation consists of:
\begin{equation}\label{eq:compressed_format}
\begin{aligned}
    W_{\mathrm{comp},i,2k+t} &= W_{i,4k+j_t},\quad t\in\{0,1\},\\
    \eta_{i,k} &= \mathrm{encode}(I_{i,k})\in\{0,\ldots,5\},
\end{aligned}
\end{equation}
where $j_0<j_1$ are the two active indices within the block, and
$\eta_{i,k}$ is a $\lceil\log_2\binom{4}{2}\rceil=3$-bit metadata selector
(in practice packed as 2-bit per element by the hardware). The resulting
compressed weight occupies $|W_{\mathrm{comp}}|=|W|/2$ storage, plus a
metadata overhead of $|\eta|=\lceil 3B/(4\times\text{element bits})\rceil$
where $B$ is the number of 4-blocks.

\paragraph{Sparse GEMM dispatch conditions.}
A layer with weight $W^{(l)}\in\mathbb{R}^{d_{\mathrm{out}}\times
d_{\mathrm{in}}}$ and mask $M^{(l)}\in\mathcal{C}_{2:4}$ is eligible for
Sparse Tensor Core dispatch when the following conditions are jointly
satisfied:
\begin{enumerate}[leftmargin=*,itemsep=1pt,parsep=0pt]
    \item \textbf{Mask validity.} $M^{(l)}\in\mathcal{C}_{2:4}$, i.e., every
          complete contiguous 4-block contains exactly two active entries.
          DST-EDL satisfies this by the top-two projection in
          Eq.~\ref{eq:mask}.
    \item \textbf{Dtype compatibility.} The weight tensor must be stored in a
          data type supported by the Sparse Tensor Core instruction set:
          $W^{(l)}\in\{\texttt{FP16},\texttt{BF16},\texttt{INT8},\texttt{FP8}\}$.
          The reported DST-EDL runs use 16-bit mixed
          precision~\cite{micikevicius2018mixed}, satisfying this requirement.
    \item \textbf{Shape alignment.} $d_{\mathrm{in}}$ must be divisible by 4
          for 2:4 block alignment. Both $d_{\mathrm{in}}$ and
          $d_{\mathrm{out}}$ should be multiples of 16 (Ampere) or 8 (Hopper)
          to fill the Tensor Core tile without padding overhead.
    \item \textbf{Layout conversion.} The dense weight $W^{(l)}$ must be
          converted to the compressed format
          $(W_{\mathrm{comp}}^{(l)},\eta^{(l)})$ before kernel dispatch. This
          is a deterministic post-training operation that preserves the
          numerical values of active weights exactly.
\end{enumerate}

\paragraph{Deployment pipeline.}
Translating a trained DST-EDL model into a hardware-accelerated inference
engine follows a three-stage pipeline:
\begin{enumerate}[leftmargin=*,itemsep=1pt,parsep=0pt]
    \item \textbf{Freeze.} After the final training epoch $T$, freeze the
          mask $M^{(T)}$ and compute the deployment weights
          $W_{\mathrm{deploy}}^{(l)}=W^{(T,l)}\odot M^{(T,l)}$ for each
          eligible layer $l$.
    \item \textbf{Compress.} Convert each
          $W_{\mathrm{deploy}}^{(l)}$ to the compressed format using either:
          \begin{itemize}[leftmargin=*,itemsep=0pt]
              \item cuSPARSELt:
                    \texttt{cusparseLtSpMMAPrune} followed by
                    \texttt{cusparseLtSpMMACompress}~\cite{nvidiacusparselt}, or
              \item PyTorch:
                    \texttt{to\_sparse\_semi\_structured($W_{\mathrm{deploy}}^{(l)}$)}~\cite{pytorchsemistructured}.
          \end{itemize}
    \item \textbf{Dispatch.} During inference, each eligible layer's
          dense GEMM $y=W_{\mathrm{deploy}}^{(l)}x+b$ is replaced by a
          sparse-dense GEMM
          $y=\mathrm{SpMM}(W_{\mathrm{comp}}^{(l)},\eta^{(l)},x)+b$
          that reads the compressed operand and uses the metadata to select
          corresponding activation elements from $x$.
\end{enumerate}
For TensorRT deployment, the frozen model is exported to ONNX and the
TensorRT engine is built with the \texttt{--sparsity=force} flag, which
instructs the builder to use Sparse Tensor Core kernels for layers with
conforming 2:4 masks.

\paragraph{Theoretical throughput bound.}
Under ideal conditions (compute-bound regime, amortized metadata overhead,
and no memory bandwidth bottleneck for the compressed operand), the Sparse
Tensor Core throughput relationship is:
\begin{equation}\label{eq:sparse_throughput}
    \mathrm{Throughput}_{\mathrm{sparse}}
    \leq
    2\times\mathrm{Throughput}_{\mathrm{dense}}.
\end{equation}
This 2$\times$ upper bound follows from the hardware design: each Sparse
Tensor Core cycle processes a compressed tile that encodes twice as many
logical elements as a dense tile of the same physical
size~\cite{mishra2021accelerating}. In practice, achievable speedup depends
on the tile fill efficiency, which varies with matrix dimensions and the
specific GPU microarchitecture~\cite{hu2024accelerating}.

\paragraph{Convolutional layer mapping.}
The cuSPARSELt library and PyTorch semi-structured kernels natively target
\texttt{nn.Linear} (GEMM) operations. For convolutional layers in the
reported ResNet-18 backbone, two mapping strategies exist:
\begin{enumerate}[leftmargin=*,itemsep=1pt,parsep=0pt]
    \item \textbf{Explicit im2col lowering.} Unfold the input feature map
          into a Toeplitz matrix and reshape the convolutional kernel
          $W\in\mathbb{R}^{C_{\mathrm{out}}\times C_{\mathrm{in}}\times
          k_h\times k_w}$ into a 2D matrix
          $\hat{W}\in\mathbb{R}^{C_{\mathrm{out}}\times
          (C_{\mathrm{in}}\cdot k_h\cdot k_w)}$. The 2:4 constraint is
          enforced along the flattened input dimension of $\hat{W}$. DST-EDL
          already applies this reshape for mask generation (Sec.~4.1).
    \item \textbf{Implicit GEMM via cuDNN.} Modern cuDNN backends can fuse
          the im2col step into the GEMM kernel. Sparse dispatch requires the
          fused kernel to accept the compressed format, which is
          backend-version dependent.
\end{enumerate}
Not all convolution configurations map efficiently to supported sparse
kernels. Small spatial dimensions ($H,W<7$), non-standard dilation, or group
convolutions may fall back to dense execution. DST-EDL keeps the first input
convolution dense, mitigating the most common failure mode (single-channel or
three-channel input layers with non-standard spatial size).

\paragraph{Scope of acceleration claims.}
The relationship between DST-EDL's structural guarantee and realized
inference speedup is summarized as follows:
\begin{itemize}[leftmargin=*,itemsep=1pt,parsep=0pt]
    \item DST-EDL guarantees $M\in\mathcal{C}_{2:4}$ by construction
          (Appendix~\ref{app:24_feasibility}), satisfying dispatch
          condition~1.
    \item The reported 16-bit mixed precision training satisfies dispatch
          condition~2.
    \item Shape alignment (condition~3) depends on the backbone architecture
          and must be verified per layer. The ResNet-18 wrapped layers use
          channel dimensions that are multiples of 64.
    \item Compressed format conversion (condition~4) is a deterministic
          post-training step that does not alter model weights or predictions.
    \item Realized speedup requires a compatible sparse kernel path and
          is not guaranteed by mask validity alone. The theoretical
          2$\times$ bound in Eq.~\ref{eq:sparse_throughput} is an upper limit
          that assumes ideal compute-bound conditions.
\end{itemize}
Table~\ref{tab:hardware_protocol_ladder} formalizes this separation into
structural (Level~1), masked diagnostic (Level~2), and realized kernel
(Level~3) evidence. The current implementation reports Levels~1 and~2.
Level~3 requires exporting frozen valid 2:4 weights to cuSPARSELt, TensorRT,
or PyTorch semi-structured backends with per-layer dispatch verification.
